\documentclass[11pt,letterpaper]{article}
\usepackage[T1]{fontenc}
\usepackage[utf8]{inputenc}
\usepackage{amsmath,amsthm,mathtools}
\usepackage{newtxtext,newtxmath}
\usepackage[margin=1in,headheight=15pt]{geometry}
\usepackage{microtype}
\usepackage{booktabs,array,longtable}
\usepackage{xcolor}
\usepackage{enumitem}
\usepackage{fancyhdr}
\usepackage{titlesec}
\usepackage[most]{tcolorbox}
\usepackage{hyperref}
\definecolor{navy}{HTML}{17324D}
\definecolor{teal}{HTML}{176B73}
\definecolor{pale}{HTML}{F0F5F7}
\hypersetup{colorlinks=true,linkcolor=navy,citecolor=teal,urlcolor=teal,
  pdftitle={Esthetic Numbers, Folded Paths, and Algebraic Diagonals},
  pdfauthor={ChatGPT},pdfsubject={Proofs of OEIS-listed A377000 conjectures and fixed-offset diagonal formulas}}
\titleformat{\section}{\large\bfseries\color{navy}}{\thesection}{0.7em}{}
\titleformat{\subsection}{\normalsize\bfseries\color{teal}}{\thesubsection}{0.6em}{}
\pagestyle{fancy}
\fancyhf{}
\fancyhead[L]{\small\scshape Esthetic numbers and algebraic diagonals}
\fancyhead[R]{\small Research note}
\fancyfoot[C]{\thepage}
\renewcommand{\headrulewidth}{0.3pt}
\setlength{\parskip}{4pt plus 1pt}
\setlength{\parindent}{1.15em}
\setlength{\emergencystretch}{2em}
\setlist{itemsep=3pt,topsep=4pt}
\numberwithin{equation}{section}
\newtheorem{theorem}{Theorem}[section]
\newtheorem{proposition}[theorem]{Proposition}
\newtheorem{lemma}[theorem]{Lemma}
\newtheorem{corollary}[theorem]{Corollary}
\theoremstyle{definition}
\newtheorem{definition}[theorem]{Definition}
\newtheorem{example}[theorem]{Example}
\newtheorem{remark}[theorem]{Remark}
\newcommand{\Z}{\mathbb Z}
\newcommand{\Q}{\mathbb Q}
\newcommand{\R}{\mathbb R}
\newcommand{\one}{\mathbf 1}
\newcommand{\coef}[2]{[#1^{#2}]}
\newcommand{\pos}[1]{(#1)_{+}}
\newcommand{\Cat}{\operatorname{Cat}}
\newcommand{\OEIS}[1]{\href{https://oeis.org/#1}{\texttt{#1}}}
\newtcolorbox{summarybox}{colback=pale,colframe=navy,boxrule=0.5pt,
  arc=1.5pt,left=10pt,right=10pt,top=8pt,bottom=8pt,breakable}

\begin{document}
\begin{titlepage}
\thispagestyle{empty}
\vspace*{0.35in}
{\color{teal}\sffamily\bfseries RESEARCH IN INTEGER SEQUENCES}\par
\vspace{0.28in}
{\fontsize{29}{34}\selectfont\bfseries\color{navy}
Esthetic Numbers,\\Folded Paths, and\\Algebraic Diagonals\par}
\vspace{0.28in}
{\Large Proofs of the listed A377000 conjectures\\and a uniform boundary-correction theorem\par}
\vspace{0.35in}
{\large ChatGPT}\par
{19 September 2026}\par
\vspace{0.35in}
\begin{abstract}
Let $T(q,k)$ count the positive base-$q$ integers with $k$ digits for which
adjacent digits differ by one.  The retrieved OEIS entry A377000 lists
conjectures relating this array to differential compositions, linear
recurrences, and the sequence A182555.  We give self-contained proofs of
these statements, with the diagonal identity as the principal target.
The key bijection identifies digit words with reflection orbits of
$k$-step walks on a path with $q$ vertices.

The same viewpoint yields a broader theorem.  On every fixed-offset
diagonal $q=k+d$, the generating function belongs to
$\Q(z,\sqrt{1-4z^2})$ and has algebraic degree exactly two.  All these
generating functions have the same irrational part, so their pairwise
differences are rational.  Below the diagonal, the correction to a
central-binomial formula is eventually a polynomial of an explicit degree;
we give its coefficients, the complete finite exceptional set, and a sharp
threshold.  We also identify the exact minimal recurrence of every row
using Chebyshev polynomials and prove the suggested A206603 identification.

A standard-library Python verifier checks the results with exact
arithmetic, including 75,381 equalities or nonsingularity checks and an
exhaustive enumeration of 131,070 sign words.  An additional symbolic
verifier checks the resolvent and reduced denominators.  The proofs, not
the experiments, establish the mathematical statements.
\end{abstract}
\vfill
\begin{summarybox}
\textbf{Status and attribution.}
This note proves the conjectures \emph{as stated in the retrieved OEIS
entry}; it does not certify worldwide priority.  The classical spectral
formula is due to De Koninck and Doyon, and the differential-operation
model is due to Male\v sevi\'c.  In particular, mere existence of row
recurrences was already an immediate consequence of known results.
The manuscript has not undergone independent peer review or
proof-assistant verification.
\end{summarybox}
\end{titlepage}

\begingroup
\small
\setlength{\parskip}{0pt}
\tableofcontents
\endgroup
\newpage

\section{The problem, the source, and the result}
\label{sec:problem}

\subsection{An explicit OEIS-listed target}
For integers $q\ge2$ and $k\ge1$, define
\begin{equation}\label{eq:definition}
T(q,k)=\#\{(d_1,\ldots,d_k):
 1\le d_1\le q-1,\quad 0\le d_i\le q-1,\quad
 |d_{i+1}-d_i|=1\}.
\end{equation}
These are the $q$-esthetic numbers of length $k$.  The initial digit is
nonzero; subsequent zero digits are allowed.  Thus $T(q,1)=q-1$, and
$T(2,k)=1$ for all positive $k$.

The array is recorded in \OEIS{A377000}, introduced by Paolo Xausa, with
three conjectures attributed to Chai Wah Wu on 21 October 2024
\cite{oeis377000}.  In the notation of this note, they assert:
\begin{enumerate}[label=\textbf{C\arabic*.},leftmargin=3em]
\item For even $q$, $T(q,k)$ counts meaningful differential compositions of
order $k$ on $\R^{q-1}$.
\item For every fixed $q$, the sequence $T(q,k)$ satisfies a linear
recurrence; several particular recurrences are displayed.
\item $T(q,q-2)=\OEIS{A182555}(q-2)$, on the domain $q\ge3$ where the
length is positive.
\end{enumerate}
The same entry suggests an identification of another diagonal with
\OEIS{A206603}.  We make its indexing precise and prove it as well.

The main chosen problem is \textbf{C3}.  It is a sharply stated equality
between a finite combinatorial count and a sequence specified by an
algebraic generating function.  Proving only a large initial segment
would not settle it; the goal is an identity of formal power series.

\subsection{What the literature check does and does not establish}
De Koninck and Doyon already express $T(q,k)$ as a finite sum of powers of
path eigenvalues \cite[Eq.~(17)]{dek}.  Consequently, C2's recurrence
\emph{existence} follows at once from prior work, even though the OEIS
entry labels it a conjecture.  We therefore do not present that part as a
new discovery.  Our recurrence treatment makes the minimal polynomial
and the parity-dependent reduction explicit.

Male\v sevi\'c's paper defines meaningful differential compositions by
matching input and output spaces and gives their adjacency rule
\cite{mal}.  We use that established definition and supply the bridge to
esthetic digit words.  The generating functions defining the two target
sequences are recorded in their own entries \cite{oeis182555,oeis206603}.

The sources and the revision history were consulted on 19 September
2026.  The retrieved history lists a 5 November 2025 link update as its
newest revision \cite{oeishistory}; the retrieved entry still labels C1--C3
as conjectures.  No separate resolution was located in the targeted
search.  This is evidence for the status of the \emph{listed problem},
not a proof that an equivalent result has never appeared elsewhere.
Historical priority of the more general diagonal formulas would require
an additional specialist literature review.

\subsection{The main extension in one statement}
Here is the structural result that goes beyond the originally selected
diagonal.  All generating functions in this note are formal series at
zero, and square roots have constant term one.

\begin{theorem}[Fixed-offset diagonals]\label{thm:overview}
For every integer $d$, put
\[
G_d(z)=\sum_{k\ge\max(1,2-d)}T(k+d,k)z^k.
\]
There is an explicitly computable rational function $R_d(z)\in\Q(z)$ such
that
\begin{equation}\label{eq:commonpart}
G_d(z)=R_d(z)-\frac{\sqrt{1-4z^2}}{2(1-2z)^2}.
\end{equation}
Thus $G_d$ has algebraic degree exactly two over $\Q(z)$, and
$G_d-G_e\in\Q(z)$ for every pair of integers $d,e$.

When $d\ge0$, one has the particularly simple formula
\begin{equation}\label{eq:positiveG}
G_d(z)=\frac{d+1-2dz-\sqrt{1-4z^2}}{2(1-2z)^2}-\frac d2.
\end{equation}
For $d=-c<0$, the rational correction is given completely in
Theorem~\ref{thm:allgf}.  Its infinite tail comes from a polynomial in $k$
of degree $\lfloor(c-1)/2\rfloor$.
\end{theorem}

The remainder of the paper proves this theorem and the OEIS
identifications.  The path bijection comes first, then the recurrence
and differential-composition consequences, followed by the diagonal
analysis.  Readers interested only in C3 may read
Sections~\ref{sec:orbits} and~\ref{sec:range}.

\section{A reflection-orbit bijection}
\label{sec:orbits}

Let $P_q$ be the path with vertices $0,1,\ldots,q-1$, and let $A_q$ be its
adjacency matrix.  Each undirected edge represents the two possible
oriented transitions between neighboring vertices.  The vector $\one$
has $q$ entries, all equal to one.

\begin{theorem}[The extra-step identity]\label{thm:orbit}
For all $q\ge2$ and $k\ge1$,
\begin{equation}\label{eq:orbit}
T(q,k)=\frac12\one^{\mathsf T}A_q^k\one.
\end{equation}
More precisely, esthetic words of length $k$ are in bijection with
reflection orbits of $k$-step walks on $P_q$.
\end{theorem}

\begin{proof}
A $k$-step walk is a sequence $(w_0,\ldots,w_k)$ of vertices whose
consecutive differences have absolute value one.  Reflect every vertex
by the map $i\mapsto q-1-i$.  This reverses the sign of the first step,
so, because $k\ge1$, no walk is fixed by reflection.  Each orbit contains
exactly one walk whose first step is upward.

Take this representative and delete $w_0$.  Since $w_1=w_0+1$, the first
remaining vertex satisfies $w_1\ge1$.  The sequence
$(w_1,\ldots,w_k)$ therefore satisfies \eqref{eq:definition}.
Conversely, prepend $d_1-1$ to any word $(d_1,\ldots,d_k)$ counted by
$T(q,k)$.  The prepended vertex lies in $\{0,\ldots,q-2\}$, and the first
step is upward.  These constructions are inverse.

Finally, $\one^{\mathsf T}A_q^k\one$ counts all $k$-step walks, with all
starting and ending vertices allowed.  Every reflection orbit has size
two, proving the formula.
\end{proof}

The exponent $k$ in \eqref{eq:orbit} is important.  A length-$k$ digit
word normally describes $k-1$ steps, with a restricted starting vertex.
The bijection replaces that asymmetry by \emph{one additional step} and
unrestricted endpoints.  It is this change of viewpoint that makes the
diagonal formulas simple.

For generating functions only, we will use the extension
\begin{equation}\label{eq:extension}
\widetilde T(q,0)=q/2,\qquad
\widetilde T(q,k)=T(q,k)\quad(k\ge1).
\end{equation}
This is half the number of zero-step walks.  It is \emph{not} a claimed
count of empty digit strings, and it need not be an integer.

\section{Differential compositions and the first conjecture}
\label{sec:differential}

\subsection{The precise meaning of a composition}
The term ``meaningful'' refers to a formally composable word of
first-order operations: the codomain of one operation must match the
domain of the next.  It does not mean that the resulting differential
operator is nonzero, nor does it identify two words that happen to define
the same operator.  This is the convention in Male\v sevi\'c's
construction \cite{mal}.

Let $q=2m$ and $N=q-1=2m-1$.  For $0\le i\le N$, identify differential
forms of degrees $i$ and $N-i$ with a common field type
$V_{\min(i,N-i)}$.  There are $m$ types $V_0,\ldots,V_{m-1}$.
Exterior differentiation, transported through these identifications,
gives one operation from $V_{\min(i-1,N-i+1)}$ to
$V_{\min(i,N-i)}$ for each $1\le i\le N$.

Its type-transition graph has one directed edge in each direction
between successive types, and one loop at the final type:
\[
V_0\ \rightleftarrows\ V_1\ \rightleftarrows\ \cdots\
\rightleftarrows\ V_{m-1},
\qquad V_{m-1}\longrightarrow V_{m-1}.
\]
The final loop contributes \emph{one} allowed transition.  No convention
that counts an undirected loop twice is being used.

For example, in dimension three the edges are
$\operatorname{grad}:V_0\to V_1$,
$\operatorname{curl}:V_1\to V_1$, and
$\operatorname{div}:V_1\to V_0$.
There are five meaningful second-order words, including
$\operatorname{curl}\operatorname{grad}$ and
$\operatorname{div}\operatorname{curl}$ even though both vanish.

\begin{theorem}[The differential-composition identification]\label{thm:differential}
For even $q\ge4$ and every $k\ge1$, the number of meaningful words of
$k$ first-order differential operations on $\R^{q-1}$ is $T(q,k)$.
The same equality holds for $q=2$ under the natural one-dimensional
interpretation consisting of the single derivative.
\end{theorem}

\begin{proof}
Write $q=2m$.  Fold the vertex set of $P_{2m}$ by reflection, mapping
$i$ to $\min(i,2m-1-i)$.  This produces exactly the type-transition graph
just described.  Every quotient vertex has two lifts.

Choose a walk in the quotient and choose one of the two lifts of its
starting vertex.  Its lift is then forced step by step.  At an ordinary
edge there is a unique neighbor mapping to the specified next vertex;
at the terminal loop there is a unique central edge crossing to the
other lift.  Consequently, every quotient walk has exactly two lifts,
which are reflections of one another.

A word of $k$ composable operations is the same as a $k$-step quotient
walk: the successive edges specify the operations and the vertices
specify their types.  Such words are therefore in bijection with
reflection orbits of $k$-step walks on $P_{2m}$.  Apply
Theorem~\ref{thm:orbit}.

For $m=1$, the quotient is a single vertex with one loop.  There is one
word of every positive length, matching $T(2,k)=1$.
\end{proof}

This proves C1.  Notice that no calculation with derivatives is needed
once the meaning of a composable word has been fixed.  The statement is
about a finite type graph, not an analytic identity between operators.

\section{Every row: spectrum and exact minimal recurrence}
\label{sec:recurrences}

\subsection{The surviving path eigenvalues}
Define
\[
\theta_j=\frac{\pi j}{q+1},\qquad \lambda_j=2\cos\theta_j,
\qquad 1\le j\le q.
\]
The vectors with components
\[
v_j(i)=\sqrt{\frac2{q+1}}\sin((i+1)\theta_j),\qquad 0\le i<q,
\]
form an orthonormal eigenbasis of $A_q$.  Indeed, the eigenvalue equation
is the sine addition formula, the endpoint values vanish at positions
$-1$ and $q$, and summing the finite geometric series gives orthogonality
and the normalization.

The finite sine sum is
\begin{equation}\label{eq:sinesum}
\sum_{i=1}^{q}\sin(i\theta_j)=
\begin{cases}
0,&j\text{ even},\\
\cot(\theta_j/2),&j\text{ odd}.
\end{cases}
\end{equation}
For completeness, start from
$\sum_{i=1}^{q}\sin(i\theta)=
\sin(q\theta/2)\sin((q+1)\theta/2)/\sin(\theta/2)$.
At $\theta=\theta_j$, the second sine vanishes for even $j$.
For odd $j$, the product in the numerator becomes
$\cos(\theta_j/2)$, giving \eqref{eq:sinesum}.

Applying the spectral decomposition to \eqref{eq:orbit} now yields
\begin{equation}\label{eq:spectral}
T(q,k)=\frac1{q+1}
\sum_{\substack{1\le j\le q\\j\text{ odd}\\\lambda_j\ne0}}
\cot^2(\theta_j/2)\lambda_j^k,\qquad k\ge1.
\end{equation}
This is the De Koninck--Doyon formula in an equivalent half-angle form
\cite[Eq.~(17)]{dek}.  We have derived it again to make the consequences
and the exceptional zero eigenvalue transparent.

\subsection{Chebyshev formulas and minimality}
Let $\mathsf T_m$ and $\mathsf U_m$ denote the Chebyshev polynomials
characterized by
\[
\mathsf T_m(\cos\theta)=\cos(m\theta),\qquad
\mathsf U_m(\cos\theta)=\frac{\sin((m+1)\theta)}{\sin\theta}.
\]
Their elementary recurrences imply that
$2\mathsf T_m(x/2)$ for $m\ge1$ and $\mathsf U_m(x/2)$ are monic integer
polynomials.

\begin{theorem}[Exact minimal row recurrence]\label{thm:minimal}
The monic minimal characteristic polynomial for the positive-length
sequence $k\mapsto T(q,k)$ is
\begin{align}
C_{2m}(x)&=\mathsf U_m(x/2)-\mathsf U_{m-1}(x/2),\label{eq:evenC}\\
C_{2m-1}(x)&=\frac{2\mathsf T_m(x/2)}{x^{m\bmod2}}.
\label{eq:oddC}
\end{align}
Its degree $d_q$ is
\begin{equation}\label{eq:degree}
d_q=\begin{cases}
q/2,&q\text{ even},\\
(q+1)/2,&q\equiv3\pmod4,\\
(q-1)/2,&q\equiv1\pmod4.
\end{cases}
\end{equation}
Writing $C_q(x)=x^{d_q}+c_1x^{d_q-1}+\cdots+c_{d_q}$, the recurrence is
\begin{equation}\label{eq:rowrecurrence}
T(q,k+d_q)+c_1T(q,k+d_q-1)+\cdots+c_{d_q}T(q,k)=0
\quad(k\ge1).
\end{equation}
No lower-order constant-coefficient recurrence holds, even eventually.
\end{theorem}

\begin{proof}
Formula~\eqref{eq:spectral} is a sum of powers of distinct, nonzero
numbers $\lambda_j$.  Each coefficient
$w_j=(q+1)^{-1}\cot^2(\theta_j/2)$ is strictly positive.  Consequently,
the product of $x-\lambda_j$ over the indicated odd indices is an
annihilating polynomial.

To prove minimality, suppose a polynomial $p$ gives a recurrence
valid for all $k\ge k_0\ge1$.  Then
\[
\sum_j w_jp(\lambda_j)\lambda_j^k=0\qquad(k\ge k_0).
\]
Use as many successive values of $k$ as there are surviving eigenvalues.
The resulting matrix is a Vandermonde matrix multiplied by a nonsingular
diagonal matrix with entries $\lambda_j^{k_0}$.  It is invertible, so
$w_jp(\lambda_j)=0$ for every $j$.  Thus $p$ vanishes at every surviving
eigenvalue.  This proves minimality of the product.

It remains to identify the product.  If $q=2m$, then
\[
\mathsf U_m(\cos\theta)-\mathsf U_{m-1}(\cos\theta)
=\frac{\cos((2m+1)\theta/2)}{\cos(\theta/2)}.
\]
Its $m$ zeros in the range under consideration occur at
$\theta=(2a-1)\pi/(2m+1)$, $1\le a\le m$.  These are exactly the odd-index
path eigenvalues, and none is zero.  Monicity proves \eqref{eq:evenC}.

If $q=2m-1$, the odd-index angles are
$(2a-1)\pi/(2m)$, $1\le a\le m$, precisely the zeros of
$\mathsf T_m(\cos\theta)$.  A zero eigenvalue is present exactly when $m$
is odd.  It contributes nothing to \eqref{eq:spectral} for $k\ge1$, so its
factor $x$ must be removed.  This gives \eqref{eq:oddC} and the three
degree cases.
\end{proof}

\begin{center}
\begin{tabular}{r l}
\toprule
$q$ & $C_q(x)$ \\
\midrule
2 & $x-1$ \\
3 & $x^2-2$ \\
4 & $x^2-x-1$ \\
5 & $x^2-3$ \\
6 & $x^3-x^2-2x+1$ \\
7 & $x^4-4x^2+2$ \\
8 & $x^4-x^3-3x^2+2x+1$ \\
9 & $x^4-5x^2+5$ \\
10 & $x^5-x^4-4x^3+3x^2+3x-1$ \\
11 & $x^6-6x^4+9x^2-2$ \\
12 & $x^6-x^5-5x^4+4x^3+6x^2-3x-1$ \\
\bottomrule
\end{tabular}
\end{center}
The displayed OEIS recurrences for $q=6,\ldots,12$ are therefore correct.
The result proves C2 while making explicit what was already implicit
in the classical spectral formula.

\begin{corollary}[Reduced denominators and a Hankel certificate]\label{cor:hankel}
The reduced denominator, normalized to have constant term one, of
$\sum_{k\ge1}T(q,k)z^k$ is $z^{d_q}C_q(1/z)$.
Moreover, for every $s\ge1$,
\begin{equation}\label{eq:hankel}
\det\bigl(T(q,s+i+j)\bigr)_{0\le i,j<d_q}
=\prod_a(w_a\lambda_a^s)\prod_{a<b}(\lambda_b-\lambda_a)^2\ne0,
\end{equation}
where the indices run through the surviving eigenvalues.
\end{corollary}
\begin{proof}
Summing $w_a\lambda_a^kz^k$ over $k\ge1$ gives
$w_a\lambda_a z/(1-\lambda_a z)$.  Its pole cannot be canceled by the
terms with other eigenvalues.  This proves the denominator claim.
The Hankel matrix factors as
$V\,\operatorname{diag}(w_a\lambda_a^s)V^{\mathsf T}$, where
$V_{ia}=\lambda_a^i$.  Taking its determinant gives \eqref{eq:hankel}.
\end{proof}

\section{Walk ranges and the diagonal conjecture}
\label{sec:range}

\subsection{Translating an unrestricted walk into a finite path}
Let $\varepsilon=(\varepsilon_1,\ldots,\varepsilon_k)\in\{-1,1\}^k$.
Set $s_0=0$, $s_i=\varepsilon_1+\cdots+\varepsilon_i$, and define
\[
L(\varepsilon)=\min_{0\le i\le k}s_i,\qquad
U(\varepsilon)=\max_{0\le i\le k}s_i,\qquad
\mathcal R(\varepsilon)=U(\varepsilon)-L(\varepsilon).
\]
Write $x_+=\max(x,0)$.

\begin{proposition}[Range formula]\label{prop:range}
For every $q\ge2$, $k\ge1$,
\begin{equation}\label{eq:range}
T(q,k)=\frac12\sum_{\varepsilon\in\{-1,1\}^k}
\bigl(q-\mathcal R(\varepsilon)\bigr)_+.
\end{equation}
If
\begin{equation}\label{eq:Mdef}
M_k=\sum_{\varepsilon\in\{-1,1\}^k}U(\varepsilon),\qquad M_0=0,
\end{equation}
then
\begin{equation}\label{eq:wide}
T(q,k)=q2^{k-1}-M_k\qquad(q\ge k,\ q\ge2).
\end{equation}
\end{proposition}
\begin{proof}
A translation by an integer $a$ fits the unrestricted walk into
$\{0,\ldots,q-1\}$ exactly when
$-L(\varepsilon)\le a\le q-1-U(\varepsilon)$.
The number of such integers is $(q-\mathcal R(\varepsilon))_+$.
Summing counts all walks on $P_q$; divide by two using
Theorem~\ref{thm:orbit}.

Since the range of a $k$-step walk is at most $k$, the positive-part
operation can be removed for $q\ge k$.  Sign reversal interchanges $U$
and $-L$, so $\sum_\varepsilon\mathcal R(\varepsilon)=2M_k$.
Substitution proves \eqref{eq:wide}.
\end{proof}

This explains why a varying base is especially tractable near and above
$q=k$.  The width constraint has no nonlinear truncation there: only a
single unrestricted-walk statistic remains.

\subsection{The total maximum of binary walks}
Introduce the formal series
\begin{equation}\label{eq:r}
r=r(z)=\frac{1-\sqrt{1-4z^2}}{2z}
=z+z^3+2z^5+5z^7+\cdots.
\end{equation}
It is the unique series with zero constant term satisfying
\begin{equation}\label{eq:rparam}
r=z(1+r^2),\qquad z=\frac r{1+r^2}.
\end{equation}

\begin{lemma}[Maximum generating function]\label{lem:maximum}
One has
\begin{equation}\label{eq:Mg}
\sum_{k\ge0}M_kz^k=\frac{r}{(1-2z)(1-r)}.
\end{equation}
Define $h_k=M_k+2^{k-1}$, including $2^{-1}=1/2$ at $k=0$.  Then
\begin{equation}\label{eq:h}
h_k=\begin{cases}
(k+\tfrac12)\binom{k}{k/2},&k\text{ even},\\[2pt]
(k+1)\binom{k}{(k-1)/2},&k\text{ odd}.
\end{cases}
\end{equation}
\end{lemma}
\begin{proof}
The generating function for walks that first hit level $1$ at their
last step is $r(z)$.  One way to see this is to remove the final upward
step and reverse the signs: the remainder is a nonnegative excursion
of even length.  Such excursions satisfy the first-return decomposition
$C(t)=1+tC(t)^2$; hence $r(z)=zC(z^2)$ and
$r=z(1+r^2)$.

A walk that first reaches level $a\ge1$ decomposes uniquely into $a$
successive first-passage segments, each climbing one new level.  Its
generating function is $r^a$.  After that first visit, an unrestricted
suffix has generating function $1/(1-2z)$.  Therefore the generating
function for walks whose maximum is at least $a$ is
$r^a/(1-2z)$.

For a nonnegative integer $U$, the identity
$U=\sum_{a\ge1}\mathbf1_{\{U\ge a\}}$ holds.  Sum over walks, then sum the
geometric series in $a$.  All sums are locally finite as formal series,
and the result is \eqref{eq:Mg}.

Using \eqref{eq:rparam} to simplify now gives
\begin{equation}\label{eq:hgf}
\sum_{k\ge0}h_kz^k
=\frac{\sqrt{1-4z^2}}{2(1-2z)^2}
=\frac{(1+2z)^2}{2(1-4z^2)^{3/2}}.
\end{equation}
The elementary expansion
\[
(1-4t)^{-3/2}=\sum_{n\ge0}(2n+1)\binom{2n}{n}t^n
\]
follows by differentiating $(1-4t)^{-1/2}$, or directly from the binomial
series.  Extracting odd powers in \eqref{eq:hgf} gives
$h_{2n+1}=2(2n+1)\binom{2n}{n}$.
Extracting even powers gives
\[
h_{2n}=\tfrac12(2n+1)\binom{2n}{n}
       +2(2n-1)\binom{2n-2}{n-1}
       =(2n+\tfrac12)\binom{2n}{n}
\]
for $n\ge1$, while $h_0=1/2$.  These are precisely \eqref{eq:h}.
\end{proof}

\subsection{A whole family above the diagonal}
For every integer $d$ and $k\ge0$, define the algebraic baseline
\begin{equation}\label{eq:baseline}
B_d(k)=(k+d+1)2^{k-1}-h_k.
\end{equation}
In particular, $B_d(0)=d/2$.  For negative $d$ or invalid small bases,
$B_d(k)$ is only an auxiliary expression, not a digit count.

\begin{proposition}[Baseline generating function]\label{prop:baseline}
For every integer $d$,
\begin{equation}\label{eq:Fd}
F_d(z):=\sum_{k\ge0}B_d(k)z^k
=\frac{d+1-2dz-\sqrt{1-4z^2}}{2(1-2z)^2}.
\end{equation}
For $d\ge0$ and valid positive lengths,
$T(k+d,k)=B_d(k)$.
\end{proposition}
\begin{proof}
The coefficient formula follows from
\eqref{eq:wide} and $M_k=h_k-2^{k-1}$.
For the generating function, sum the linear-exponential term in
\eqref{eq:baseline} and subtract \eqref{eq:hgf}:
\[
\sum_{k\ge0}(k+d+1)2^{k-1}z^k
=\frac{z}{(1-2z)^2}+\frac{d+1}{2(1-2z)}.
\]
This simplifies to \eqref{eq:Fd}.
\end{proof}

\begin{theorem}[The selected diagonal conjecture]\label{thm:C3}
For every $k\ge1$,
\begin{equation}\label{eq:C3}
T(k+2,k)=\OEIS{A182555}(k).
\end{equation}
Equivalently, C3 holds for all $q\ge3$.
\end{theorem}
\begin{proof}
Set $d=2$ in Proposition~\ref{prop:baseline}.  Since $B_2(0)=1$,
\begin{equation}\label{eq:A182GF}
1+\sum_{k\ge1}T(k+2,k)z^k
=\frac{3-4z-\sqrt{1-4z^2}}{2(1-2z)^2}.
\end{equation}
The right-hand side is the defining generating function of
\OEIS{A182555} \cite{oeis182555}.  Equality of formal power series proves
equality of every coefficient, not merely a finite initial segment.
\end{proof}

An explicit positive formula is also available.  Let
$\Cat(j)=\binom{2j}{j}/(j+1)$.  Since
$(1-\sqrt{1-4z^2})/2=\sum_{j\ge1}\Cat(j-1)z^{2j}$,
\begin{equation}\label{eq:positivecatalan}
T(k+2,k)=2^k+
\sum_{j=1}^{\lfloor k/2\rfloor}
\Cat(j-1)(k-2j+1)2^{k-2j}.
\end{equation}
This gives a finite-sum expression independent of the original digit
recursion.

\subsection{The suggested A206603 diagonal}
\begin{theorem}[An addition-triangle diagonal]\label{thm:A206603}
For $k\ge2$,
\begin{equation}\label{eq:A206}
T(k,k)=\OEIS{A206603}(k),\qquad
\sum_{k\ge2}T(k,k)z^k
=\frac{1-\sqrt{1-4z^2}}{2(1-2z)^2}.
\end{equation}
Moreover,
\begin{equation}\label{eq:diagdiff}
T(k+2,k)-T(k,k)=2^k\qquad(k\ge2).
\end{equation}
\end{theorem}
\begin{proof}
Set $d=0$ in \eqref{eq:Fd}.  Its constant and linear coefficients vanish,
so it is already the generating function of the valid diagonal.
It matches the generating function recorded for \OEIS{A206603}
\cite{oeis206603}.  Subtract the two baseline formulas with $d=2$ and
$d=0$ to obtain \eqref{eq:diagdiff}.
\end{proof}

There is also a direct check against the \emph{combinatorial definition}
of A206603, without taking its generating function on trust.  In an
addition triangle with $n+1$ entries in the base, the apex weights those
entries by $\binom n0,\ldots,\binom nn$.  Pairwise swaps show that the
maximum is attained by assigning the increasing base values
$j-n/2$ to the increasing binomial weights.  The sorted weights are
$\binom n{\lfloor j/2\rfloor}$, $0\le j\le n$.

For $n=2m$, pair equal weights to obtain
\[
\sum_{j=0}^{m-1}(4j+1-2m)\binom{2m}{j}
+m\binom{2m}{m}
=(2m+1)2^{2m-1}-(2m+\tfrac12)\binom{2m}{m}.
\]
For $n=2m+1$, the corresponding paired sum is
\[
\sum_{j=0}^{m}(4j-2m)\binom{2m+1}{j}
=(2m+2)2^{2m}-2(2m+1)\binom{2m}{m}.
\]
These identities follow from binomial symmetry and
$j\binom nj=n\binom{n-1}{j-1}$.  The two results are exactly $B_0(n)$
from \eqref{eq:h}--\eqref{eq:baseline}.  Thus the maximum apex and the
esthetic diagonal agree directly as well.

\section{An exact finite-width resolvent and reflection formula}
\label{sec:resolvent}

For $q<k$, the positive-part operation in \eqref{eq:range} matters.
We now compute its effect exactly.  This provides the generalization
needed for negative-offset diagonals.

\subsection{Solving the endpoint-value recurrence}
For vertices indexed by $1,\ldots,q$, let
\[
v_i(z)=\sum_{k\ge0}(A_q^k\one)_iz^k.
\]
Then
\begin{equation}\label{eq:vrec}
v_i=1+z(v_{i-1}+v_{i+1}),\qquad v_0=v_{q+1}=0.
\end{equation}
The equation has a unique solution in formal power series: its constant
terms are fixed, and each successive coefficient is determined by the
previous coefficients.

\begin{proposition}[Exact row resolvent]\label{prop:resolvent}
With $r=r(z)$ as in \eqref{eq:r},
\begin{equation}\label{eq:visol}
v_i(z)=\frac{1-(r^i+r^{q+1-i})/(1+r^{q+1})}{1-2z}.
\end{equation}
Consequently,
\begin{align}
\mathcal T_q(z)
&:=\frac q2+\sum_{k\ge1}T(q,k)z^k\notag\\
&=\frac{q}{2(1-2z)}
-\frac{r(1-r^q)}{(1-2z)(1-r)(1+r^{q+1})}.
\label{eq:rowgf}
\end{align}
\end{proposition}
\begin{proof}
The right-hand side of \eqref{eq:visol} vanishes at the two boundary
indices.  Its constant particular solution is $1/(1-2z)$.
The remaining terms satisfy the homogeneous recurrence because
$r^i=z(r^{i-1}+r^{i+1})$, and the same is true of $r^{q+1-i}$.
Thus it solves \eqref{eq:vrec}, and uniqueness proves the assertion.
Sum over $i$ and divide by two, using
$\sum_{i=1}^q r^i=r(1-r^q)/(1-r)$ and
Theorem~\ref{thm:orbit}.
\end{proof}

Although \eqref{eq:rowgf} is expressed using an algebraic series $r$, the
whole expression is rational in $z$ for fixed $q$, as is already clear
from the finite matrix resolvent.  It should not be confused with the
genuinely nonrational \emph{diagonal} generating functions.

Define
\begin{equation}\label{eq:K}
K(u)=\frac{(1+u)(1+u^2)}{(1-u)^3}.
\end{equation}
Using $1-2z=(1-r)^2/(1+r^2)$, rewrite the row resolvent as
\begin{equation}\label{eq:rowsplit}
\mathcal T_q(z)
=\frac{q}{2(1-2z)}-\frac{r}{(1-2z)(1-r)}
+K(r)\frac{r^{q+1}}{1+r^{q+1}}.
\end{equation}
The first two terms have coefficient $q2^{k-1}-M_k$.  The third term is
therefore precisely the finite-width correction.

\subsection{A coefficient-extraction identity}
Set
\begin{equation}\label{eq:beta}
\beta_0=1,\qquad \beta_t=4t\quad(t\ge1),
\end{equation}
and for $k\ge0$ define
\begin{equation}\label{eq:S}
S(k,e)=\begin{cases}
\displaystyle\sum_{\ell=0}^{\lfloor e/2\rfloor}
\beta_{e-2\ell}\binom{k}{\ell},&e\ge0,\\[6pt]
0,&e<0.
\end{cases}
\end{equation}
As usual, $\binom{k}{\ell}=0$ when $\ell>k\ge0$.

\begin{lemma}[The correction kernel]\label{lem:kernel}
For integers $a\ge1$ and $k\ge0$,
\begin{equation}\label{eq:kernel}
\coef{z}{k}K(r(z))r(z)^a
=\coef{u}{k-a}\frac{(1+u)^2}{(1-u)^2}(1+u^2)^k
=S(k,k-a).
\end{equation}
\end{lemma}
\begin{proof}
The change of variables $z=u/(1+u^2)$ is valid for formal residues,
because its linear coefficient is one.  Its differential is
$dz=(1-u^2)(1+u^2)^{-2}\,du$.  Therefore
\begin{align*}
\coef{z}{k}K(r(z))r(z)^a
&=\operatorname{Res}_{u=0}
K(u)u^a\frac{(1+u^2)^{k+1}}{u^{k+1}}
\frac{1-u^2}{(1+u^2)^2}\,du\\
&=\operatorname{Res}_{u=0}
u^{a-k-1}\frac{(1+u)^2}{(1-u)^2}(1+u^2)^k\,du.
\end{align*}
This proves the first equality.  Finally,
\[
\frac{(1+u)^2}{(1-u)^2}=1+4\sum_{t\ge1}tu^t
=\sum_{t\ge0}\beta_tu^t.
\]
Multiplication by the binomial expansion of $(1+u^2)^k$ proves the second.
\end{proof}

\begin{theorem}[Universal finite reflection sum]\label{thm:reflection}
For every $q\ge2$ and $k\ge1$,
\begin{equation}\label{eq:universal}
\boxed{\displaystyle
T(q,k)=q2^{k-1}-M_k+
\sum_{j=1}^{\lfloor k/(q+1)\rfloor}
(-1)^{j-1}S\bigl(k,k-j(q+1)\bigr).}
\end{equation}
Equivalently, replace $q2^{k-1}-M_k$ by $(q+1)2^{k-1}-h_k$.
\end{theorem}
\begin{proof}
Expand the last term of \eqref{eq:rowsplit} as
\[
K(r)\frac{r^{q+1}}{1+r^{q+1}}
=\sum_{j\ge1}(-1)^{j-1}K(r)r^{j(q+1)}.
\]
This is a well-defined formal series since $r$ has order one.  For a
coefficient of degree $k$, only $j(q+1)\le k$ can contribute.
Apply Lemma~\ref{lem:kernel} with $a=j(q+1)$ and use
Lemma~\ref{lem:maximum} for the remaining terms.
\end{proof}

The formula is an identity for all widths, not an asymptotic expansion.
In particular, the apparently alternating correction is always equal
to the nonnegative quantity
\begin{equation}\label{eq:overflow}
\frac12\sum_{\varepsilon\in\{-1,1\}^k}
\bigl(\mathcal R(\varepsilon)-q\bigr)_+,
\end{equation}
because $(q-R)_+=q-R+(R-q)_+$.  It compensates for the negative
contributions that the untruncated baseline would otherwise assign to
walks of excessive range.

\section{Polynomial boundary corrections and all diagonal generating functions}
\label{sec:diagonals}

\subsection{The eventual correction polynomial}
For $c\ge1$, define a polynomial in $x$ by
\begin{equation}\label{eq:Pc}
P_c(x)=\sum_{\ell=0}^{\lfloor(c-1)/2\rfloor}
\beta_{c-1-2\ell}\binom{x}{\ell},
\qquad
\binom{x}{\ell}=\frac{x(x-1)\cdots(x-\ell+1)}{\ell!}.
\end{equation}
Thus $P_c(k)=S(k,c-1)$ for nonnegative integers $k$.

\begin{theorem}[Polynomial correction with a sharp threshold]\label{thm:polynomial}
For $c\ge1$ and
\begin{equation}\label{eq:threshold}
k\ge K_c:=\max(c+2,2c-1),
\end{equation}
one has
\begin{equation}\label{eq:polynomial}
\boxed{\displaystyle
T(k-c,k)=B_{-c}(k)+P_c(k)
=(k-c+1)2^{k-1}-h_k+P_c(k).}
\end{equation}
The degree of $P_c$ is $\lfloor(c-1)/2\rfloor$.
For $c\ge4$, the threshold is sharp in the following strong sense:
\begin{equation}\label{eq:sharp}
T(c-2,2c-2)=B_{-c}(2c-2)+P_c(2c-2)-1.
\end{equation}
\end{theorem}
\begin{proof}
Put $q=k-c$ in \eqref{eq:universal}.  The requirement $q\ge2$ is
$k\ge c+2$.  The $j=1$ term has coefficient index
$k-(q+1)=c-1$, so it is $P_c(k)$.
All terms with $j\ge2$ vanish if
$2(q+1)>k$, equivalently $k>2c-2$.  This proves the formula on exactly
the range \eqref{eq:threshold}.

Write $p=\lfloor(c-1)/2\rfloor$.  The coefficient of
$\binom{x}{p}$ is $\beta_0=1$ for odd $c$, and $\beta_1=4$ for even $c$.
It is nonzero; the degree is therefore $p$, with leading coefficient
$1/p!$ or $4/p!$, respectively.

For $c\ge4$ take $k=2c-2$, so $q=c-2\ge2$.
Then $k=2(q+1)$.  Besides the $j=1$ term there is exactly one further
term, namely the $j=2$ contribution $-S(k,0)=-1$.
No higher term can occur.  This proves \eqref{eq:sharp}.
\end{proof}

\begin{center}
\begin{tabular}{r l r}
\toprule
$c$ & $P_c(k)$ & First guaranteed valid $k=K_c$ \\
\midrule
1 & $1$ & 3 \\
2 & $4$ & 4 \\
3 & $k+8$ & 5 \\
4 & $4k+12$ & 7 \\
5 & $\binom{k}{2}+8k+16$ & 9 \\
6 & $4\binom{k}{2}+12k+20$ & 11 \\
7 & $\binom{k}{3}+8\binom{k}{2}+16k+24$ & 13 \\
8 & $4\binom{k}{3}+12\binom{k}{2}+20k+28$ & 15 \\
\bottomrule
\end{tabular}
\end{center}

\begin{example}[Why the finite exception cannot be ignored]
For $c=4$, the first admissible length is $k=6$, with base $q=2$.
Here $P_4(6)=36$, while
$B_{-4}(6)=3\cdot32-(6+\tfrac12)\binom63=96-130=-34$.
Thus $B_{-4}(6)+P_4(6)=2$, whereas $T(2,6)=1$.
The missing $-1$ is exactly the second reflection term in
\eqref{eq:sharp}.  At every $k\ge7$, the polynomial formula is exact.
\end{example}

\subsection{An explicit formula including every exceptional coefficient}
For $c\ge1$ set
\begin{equation}\label{eq:Pg}
P_c^*(z)=\sum_{\ell=0}^{\lfloor(c-1)/2\rfloor}
\beta_{c-1-2\ell}\frac{z^\ell}{(1-z)^{\ell+1}}.
\end{equation}
It is the generating function of $P_c(k)$, because
$\sum_{k\ge0}\binom{k}{\ell}z^k=z^\ell/(1-z)^{\ell+1}$.
Define the finite polynomial
\begin{equation}\label{eq:Ec}
E_c(z)=\sum_{k=c+2}^{2c-2}z^k
\sum_{j=2}^{\lfloor k/(k-c+1)\rfloor}
(-1)^{j-1}S\bigl(k,k-j(k-c+1)\bigr),
\end{equation}
with any empty sum equal to zero.  In particular, $E_c=0$ for $c\le3$.

\begin{theorem}[Every fixed-offset generating function]\label{thm:allgf}
With the valid-domain convention of Theorem~\ref{thm:overview},
\begin{equation}\label{eq:Gpositive}
G_d(z)=F_d(z)-d/2\qquad(d\ge0),
\end{equation}
and, for every $c\ge1$,
\begin{equation}\label{eq:Gnegative}
\begin{split}
G_{-c}(z)={}&F_{-c}(z)+P_c^*(z)\\
&-\sum_{k=0}^{c+1}\bigl(B_{-c}(k)+P_c(k)\bigr)z^k+E_c(z).
\end{split}
\end{equation}
Consequently, all conclusions of Theorem~\ref{thm:overview} hold.
\end{theorem}
\begin{proof}
For $d\ge0$, the baseline counts every valid positive-length term.
Remove its constant term $d/2$.  The only further potentially invalid
positive length is $k=1,d=0$, but $B_0(1)=0$, so it requires no additional
subtraction.  This proves \eqref{eq:Gpositive}.

For $d=-c$, the valid lengths begin at $c+2$.  Start with the series
$F_{-c}+P_c^*$, whose coefficient of $z^k$ is
$B_{-c}(k)+P_c(k)$.  Subtract all coefficients below the valid range;
this is exactly the finite sum in \eqref{eq:Gnegative}.
At valid lengths $k\ge2c-1$, Theorem~\ref{thm:polynomial} says that no
more changes are needed.  At the finitely many remaining valid lengths,
Theorem~\ref{thm:reflection} says that the missing terms are exactly
those with $j\ge2$ recorded in $E_c$.  Thus every coefficient is correct.

All terms of these formulas other than the square root in $F_d$ are
rational.  The square-root term is independent of $d$, proving
\eqref{eq:commonpart} and rationality of every difference.
It also proves algebraicity of degree at most two.  The degree is
\emph{exactly} two: if $G_d$ were rational, then
\eqref{eq:commonpart} would make $\sqrt{1-4z^2}$ rational.
But a square of a rational function has even order at every zero or
pole, while $1-4z^2$ has simple zeros at $z=\pm1/2$.  This is impossible.
\end{proof}

These formulas give a constructive procedure for $R_d(z)$, not only an
existence statement.  Once $R_d$ is written down, a polynomial equation
for $G_d$ is
\begin{equation}\label{eq:quadratic}
\bigl(2(1-2z)^2(G_d(z)-R_d(z))\bigr)^2=1-4z^2.
\end{equation}
Multiplying by the squared denominator of $R_d$ makes it a polynomial
identity with coefficients in $\Q[z]$.

For the first three negative offsets, the exceptional reflection
polynomial $E_c$ vanishes, but the low-index subtraction must still be
performed.  Direct evaluation gives
\begin{align*}
G_{-1}(z)&=F_{-1}(z)+\frac1{1-z}-\frac12,\\
G_{-2}(z)&=F_{-2}(z)+\frac4{1-z}-3-2z-z^2,\\
G_{-3}(z)&=F_{-3}(z)+\frac{z}{(1-z)^2}+\frac8{1-z}\\
&\quad-\frac{13}{2}-6z-5z^2-3z^3-z^4.
\end{align*}
These examples also show why the artificial length-zero extension must
be separated from the combinatorial series.

\section{Consequences and computational use}
\label{sec:consequences}

\subsection{Diagonal differences satisfy constant-coefficient recurrences}
Define $Q_d(k)=0$ for $d\ge0$, and $Q_{-c}(k)=P_c(k)$ for $c\ge1$.
The preceding results give, for every fixed pair $d,e$ and all
sufficiently large $k$,
\begin{equation}\label{eq:pairdifference}
T(k+d,k)-T(k+e,k)
=(d-e)2^{k-1}+Q_d(k)-Q_e(k).
\end{equation}
Thus the nonrational component disappears not just at the generating
function level, but directly in the exact coefficient formulas.

If the polynomial on the right has degree at most $r$, the sequence of
differences is annihilated eventually by
\[
(\mathcal E-2)(\mathcal E-1)^{r+1},\qquad
(\mathcal E a)(k)=a(k+1).
\]
When the polynomial is zero, the factor $(\mathcal E-2)$ alone suffices.
For two distinct nonnegative offsets, \eqref{eq:pairdifference} is exact
at every length where both bases are valid, and reduces to
$(d-e)2^{k-1}$.

\subsection{A common diagonal asymptotic scale}
The exact central-binomial formulas show that fixed offsets share the
same leading boundary loss.  Stirling's expansion applied separately
to the two parities gives
\begin{equation}\label{eq:hasymp}
h_k=2^k\sqrt{\frac{2k}{\pi}}
\left(1+\frac1{4k}+O(k^{-2})\right).
\end{equation}
For even $k$, insert
$\binom{k}{k/2}=2^k\sqrt{2/(\pi k)}(1-1/(4k)+O(k^{-2}))$
into \eqref{eq:h}.  For odd $k$, use
$h_k=2k\binom{k-1}{(k-1)/2}$ and expand in $1/k$; the coefficient
$1/(4k)$ is the same.

Therefore, for every fixed integer $d$,
\begin{equation}\label{eq:diagasy}
T(k+d,k)=(k+d+1)2^{k-1}
-2^k\sqrt{\frac{2k}{\pi}}+O_d(2^kk^{-1/2}).
\end{equation}
For negative offsets, the polynomial $Q_d(k)$ is exponentially smaller
than the displayed error.  When exact low-order effects matter, it
should instead be retained explicitly through \eqref{eq:polynomial}.
The asymptotic estimate is a consequence of an exact formula, not the
basis of the proof of a diagonal identity.

\subsection{Which exact algorithm to use}
The digit dynamic program starts with endpoint counts
$(0,1,\ldots,1)$ and repeatedly replaces the count at each vertex by the
sum at its neighbors.  It uses $O(qk)$ integer additions and $O(q)$
stored integers to produce a length-$k$ count.  The folded type graph
halves the state count for even $q$.

For a fixed base and a very large isolated index, the minimal recurrence
can be used with polynomial reduction modulo $C_q$.  Binary powering
computes $x^k\bmod C_q(x)$ in $O(d_q^2\log k)$ arithmetic operations
using ordinary polynomial multiplication, after which the result is a
linear combination of initial terms.  This is an arithmetic-operation
bound, \emph{not} a logarithmic bit-complexity claim: the output itself
has a number of bits growing with $k$.

For fixed offsets, the central-binomial formula with the explicit
polynomial correction is usually the most direct description.  It
requires no table over all smaller bases.  Conversely, the universal
reflection sum can involve large canceling quantities when the base
is small; it is valuable as an exact formula and proof device, but is
not asserted to be the fastest numerical method in every regime.

\section{Exact verification and reproducible artifacts}
\label{sec:verification}

\subsection{Independent models and test coverage}
The file \texttt{verify.py} uses only the Python standard library.
The primary counts come from the endpoint digit recurrence, not from
the new closed forms.  Separate checks use unrestricted path walks,
the folded graph, and Male\v sevi\'c's operation-index adjacency rule
\[
M_{ij}=1\quad\Longleftrightarrow\quad
j=i+1\ \text{or}\ i+j=N+1,
\qquad 1\le i,j\le N.
\]
This last model counts words by $\one^{\mathsf T}M^{k-1}\one$ and is
independent of the type-vertex implementation.

The complete test run performed for this package passed the following
checks.  An ``equality'' refers to one exact comparison, not a
floating-point tolerance test.

\begin{center}
\small
\begin{tabular}{p{0.48\textwidth}r p{0.27\textwidth}}
\toprule
Check & Number & Range or interpretation \\
\midrule
Reflection-orbit count & 24,750 & $2\le q\le100$, $1\le k\le250$ \\
Minimal-polynomial recurrence & 22,225 & Same rows; every fitting window \\
Differential-operation graph & 2,000 & Even $2\le q\le40$, $k\le100$ \\
Folded-path graph & 2,000 & Same even bases and lengths \\
Total maximum by enumeration & 16 & All sign words, $1\le k\le16$ \\
Range formula & 320 & $2\le q\le21$, $1\le k\le16$ \\
Universal reflection formula & 5,850 & $2\le q\le40$, $1\le k\le150$ \\
Nonnegative-offset formula & 3,249 & $0\le d\le12$, valid $k\le250$ \\
Eventual polynomial formula & 6,627 & $1\le c\le30$, valid tail $k\le250$ \\
Complete negative-diagonal series & 7,530 & Includes every invalid or exceptional low index \\
Sharp threshold & 27 & $4\le c\le30$ \\
A182555 generating function & 250 & Through index 250 \\
A206603 generating function & 249 & Indices $2$ through $250$ \\
Addition-triangle maximum & 249 & Independently sorted binomial weights \\
Nonzero minimal-size Hankel determinant & 39 & Bases $2$ through $40$ \\
\midrule
\textbf{Total} & \textbf{75,381} & Exact equalities or nonsingularity checks \\
\bottomrule
\end{tabular}
\end{center}

The exhaustive part examines $\sum_{k=1}^{16}2^k=131{,}070$ sign words.
The Hankel determinants use fraction-free integer elimination.
The recorded run used Python 3.13.5; the verifier completed in about
2.7 seconds in the working environment.  Timing depends on hardware
and is not a complexity claim.

\subsection{Symbolic checks and their limits}
The optional file \texttt{symbolic\_verify.py}, run with SymPy 1.14.0,
inverts $I-zA_q$ exactly for $2\le q\le14$.  It checks the row-resolvent
formula after the rational substitution $z=u/(1+u^2)$, verifies the
reduced denominator against the Chebyshev minimal polynomial, and checks
three scalar rational identities used in the coefficient extraction.
All of these checks passed.

The verifier intentionally separates empirical validation from proof.
No finite collection of tests establishes a statement for all $q,k$;
that role belongs to the bijections, formal-series identities, and
minimality argument above.  Likewise, neither the code nor the proofs
can certify historical novelty.  No Lean or other proof-assistant
formalization is included or claimed.

\subsection{Contents of the archive}
The archive contains the editable \LaTeX\ source and compiled PDF,
the two verifier programs, a build file, a source-provenance note, and
machine-readable results.  The generated data include:
\begin{itemize}
\item row values for $2\le q\le40$ and $1\le k\le100$;
\item diagonal values for offsets $-12\le d\le12$ through length $200$;
\item all minimal polynomials for $2\le q\le100$ and the exact Hankel
  determinants used in the finite checks.
\end{itemize}
Run \texttt{python3 verify.py} from the archive directory to regenerate
the standard-library results.  The \texttt{README.md} supplies the
complete build and reproduction commands.  No external source papers
or font files are redistributed.

\section{Conclusion}
The chosen OEIS-listed diagonal conjecture has a short conceptual
resolution: append an initial upward step, pass to reflection orbits,
and translate unrestricted binary walks into a finite path.  At width
$q\ge k$, the translation count is linear in the walk range, and a
first-passage generating function evaluates the required total maximum.
This proves C3 by exact generating-function equality.

The same construction proves C1 and gives explicit minimal recurrences
for C2, with appropriate attribution to the known spectral theory.
More importantly, it extends the diagonal identity to every fixed
integer offset.  Negative offsets introduce an explicit polynomial
boundary correction after a sharp threshold, and only finitely many
additional reflection terms before it.  The resulting generating
functions all share one quadratic irrational component.

The mathematical claims in this note are supported by complete
arguments and reproducible exact checks.  Their publication history is
a separate question: the paper should be regarded as a proposed
self-contained resolution and extension of the retrieved OEIS
conjectures, not as a certificate of priority.

\appendix
\section{Compact statements suitable for an OEIS update}
The following formulas summarize the results without depending on the
auxiliary derivations.  They are a draft for review, not a claim that an
OEIS edit has been submitted.

For $q\ge2$ and $k\ge1$, let $A_q$ be the adjacency matrix of the path
on $q$ vertices.  Then
\[
T(q,k)=\tfrac12\one^{\mathsf T}A_q^k\one.
\]
For even $q$, folding the path by reflection gives the type graph of
meaningful differential compositions on $\R^{q-1}$, proving Conjecture~1.

The minimal row characteristic polynomials are
\[
C_{2m}(x)=\mathsf U_m(x/2)-\mathsf U_{m-1}(x/2),\qquad
C_{2m-1}(x)=2\mathsf T_m(x/2)/x^{m\bmod2}.
\]
They prove the listed row recurrences and specify the exact minimal
orders.  Recurrence existence already follows from the known spectral
formula.

For $q\ge k\ge1$ and $q\ge2$,
\[
T(q,k)=(q+1)2^{k-1}-
\begin{cases}
(k+\tfrac12)\binom{k}{k/2},&k\text{ even},\\
(k+1)\binom{k}{(k-1)/2},&k\text{ odd}.
\end{cases}
\]
In particular,
\[
1+\sum_{k\ge1}T(k+2,k)z^k
=\frac{3-4z-\sqrt{1-4z^2}}{2(1-2z)^2},
\]
proving $T(q,q-2)=\OEIS{A182555}(q-2)$ for $q\ge3$.
Also $T(k,k)=\OEIS{A206603}(k)$ for $k\ge2$.
All fixed-offset diagonals $q=k+d$, including negative $d$, have
algebraic generating functions of degree two, with the explicit
formulas in Theorem~\ref{thm:allgf}.

\begin{thebibliography}{9}
\addcontentsline{toc}{section}{References}
\bibitem{oeis377000}
The OEIS Foundation Inc.,
\emph{A377000: Array read by ascending antidiagonals: $T(n,k)$ = number of
$n$-esthetic numbers with $k$ digits}.
Entry by Paolo Xausa, 12 October 2024; conjectures by Chai Wah Wu,
21 October 2024.  Retrieved 19 September 2026.
\url{https://oeis.org/A377000}.

\bibitem{oeishistory}
The OEIS Foundation Inc., \emph{Revision history for A377000}.
Retrieved 19 September 2026.
\url{https://oeis.org/history?seq=A377000}.

\bibitem{dek}
Jean-Marie De Koninck and Nicolas Doyon,
\emph{Esthetic Numbers},
Annales des sciences math\'ematiques du Qu\'ebec
\textbf{33} (2009), no.~2, 155--164.
\url{https://www.labmath.uqam.ca/~annales/volumes/33-2/PDF/155-164.pdf}.

\bibitem{mal}
Branko J. Male\v sevi\'c,
\emph{Some combinatorial aspects of differential operation composition
on the space $\R^n$},
Univ. Beograd. Publ. Elektrotehn. Fak., Ser. Mat.
\textbf{9} (1998), 29--33.
ArXiv version posted in 2007: \url{https://arxiv.org/abs/0704.0750}.

\bibitem{oeis182555}
The OEIS Foundation Inc.,
\emph{A182555: Expansion of g.f.
$(3-4x-\sqrt{1-4x^2})/(2(1-2x)^2)$}.
Entry by N. J. A. Sloane, 5 May 2012.  Retrieved 19 September 2026.
\url{https://oeis.org/A182555}.

\bibitem{oeis206603}
The OEIS Foundation Inc.,
\emph{A206603: Maximal apex value of an addition triangle whose base is a
permutation of $\{k-n/2:k=0,\ldots,n\}$}.
Entry by Alois P. Heinz, 10 February 2012.  Retrieved 19 September 2026.
\url{https://oeis.org/A206603}.
\end{thebibliography}
\end{document}
